% !TeX root = ../../Book2.tex
\chapter*{序言}
\addcontentsline{toc}{chapter}{序言}
\markboth{序言}{序言}

在向量空间里，非零标量总能除掉。把标量换成整数，这件事就不再理所当然了：方程 \(2x=1\) 在整数中无解，\(2\mathbb Z\) 虽然与 \(\mathbb Z\) 作为交换群同构，却不是 \(\mathbb Z\) 的一个直和分量。我们熟悉的基、维数和补空间，也因此需要重新检查。本卷研究环作用下的模，并考察向量空间中的哪些结论仍然成立。

\ProseParagraphBreak
整数模、向量空间和带一个线性算子的向量空间，看似是三类对象。前者的标量来自 \(\mathbb Z\)，第二类来自域 \(k\)；在第三类中，若把多项式 \(f(t)\) 的作用定义为 \(f(T)\)，标量便来自 \(k[t]\)。它们都可以用模的语言讨论，却各有自己的性质。域上每个向量空间都有基；一般环上的模未必自由，即使由有限多个元素生成，也可能有不能忽略的关系。抽象定义在这里并没有抹平差别，它让差别有了可以逐项比较的形式。

\ProseParagraphBreak
十九世纪的线性代数留下了 Smith 标准形与 Jordan 标准形这两个名字。前者整理整数或多项式系数的关系矩阵，后者描述多项式分裂后的线性算子。乍看是两套算法，把算子看成 \(k[t]\) 模以后，却能从主理想整环上有限生成模的分类理解它们之间的联系。本卷先建立自由模、展示和正合列，再回到矩阵：为什么能够化简，化简所得的数据为什么不随所选生成元改变，标准形又怎样回答最初的相似问题，都须在这一往返中说清楚。

\ProseParagraphBreak
从一个变量转向两个变量，还会遇到另一种关系。设 \(R\) 是交换环，\(M,N\) 是 \(R\) 模。双线性映射每次要同时处理两个变量；张量积 \(M\otimes_RN\) 将这种要求变成关于一个模的线性映射问题。在非交换环上，左作用和右作用却不能互换，连张量积写下之前都要先核对两侧标量是否相配。初次读到这里，不妨停下来亲手检查一个平衡关系。许多以后看来顺手的公式，正是靠这种检查才成立。

\ProseParagraphBreak
同样需要追问的，还有“解得出来”是什么意思。一个模是投射的，意味着某些映射可以沿满射提升；一个模是内射的，意味着映射可以从子模向外延拓；张量积保持短正合列左端的能力，则由平坦性描述。这些条件都与映射有关，但不能因为名称并列就把它们混为一谈。书中先用具体的自由模、商模和短正合列训练这些判别，再说明各自适用的范围。第四卷会在此基础上继续研究消解及导出函子；读第二卷时，无须预先掌握那套理论。

\ProseParagraphBreak
多线性代数研究同时依赖几个向量的映射。张量代数允许依次相乘，对称代数要求某些因子能够交换，外代数则把重复方向的乘积置为零；行列式可以从最高外幂的作用中读出。这样理解以后，计算公式并没有消失，只是知道了它们为何在换基后仍成立。双线性型和二次型又提出不同的问题：改变基时，一个算子用相似变换比较，一个型用合同变换比较。两者都写成矩阵，却不能按同一套不变量分类。

\ProseParagraphBreak
进入非交换环，乘法的次序开始影响每一步。矩阵环是熟悉的例子：\(AB\) 与 \(BA\) 往往不同，左模与右模也须分清。有限维代数中，幂等元帮助我们拆开正则模，半单性让简单模真正成为直和分量；Jacobson 根则记录了半单商看不见的部分。Artin--Wedderburn 定理把半单环写成除环上的矩阵环乘积，但这个结论有自己的条件，不能拿来替所有环作答。面对一个具体代数，先找出它的简单模、根和幂等元，通常比急于背下分类式更有收获。

\ProseParagraphBreak
Morita 理论研究一种更宽的等价：两个环不必同构，模范畴却可以等价。它促使我们辨别哪些性质由模及其同态决定。中心单代数与 Brauer 群则研究底域扩张带来的变化：一个代数在原域上未必是矩阵环，扩域后却可能变成矩阵环。这些主题各有独立的问题。最后两章的非交换分式、恒等式和滤过方法供选读；对某项计算真正有需要时，再沿着相应章节深入即可。

\ProseParagraphBreak
本卷面向学过本科线性代数、能读基本证明的读者。开始第一编之前，应熟悉第一卷中的群与环、理想和商环、整除及一元多项式；学习第四编的 Brauer 群时，再补上有限 Galois 扩张、范数等知识。模的定义从头建立，因此不要求先读完第一卷的高级域论。范畴语言在第五章由模的实例引入，也无须预先阅读第四卷。各编导读会说明所需基础和可以暂缓阅读的内容。

\ProseParagraphBreak
想先获得一套可动手使用的工具，可以顺读第一至第七章，再读第八至第十章的多线性构造和型。第一章重温线性代数，第二、三章建立模及其关系，第四章把这些工具用于标准形；第五至七章讨论自然性、张量积与正合性。已有较扎实线性代数基础的读者，可以在第一章挑出商空间、对偶和算子模作复习，然后从第二章开始。第三编的半单性与根，适合在这些基础上继续；第四编的 Morita 理论和中心单代数则可按兴趣分别进入，不必为读其中一章先通读另一章。

\ProseParagraphBreak
读模论时，不妨反复问三个问题：标量来自哪个环，作用在左边还是右边，所写的元素有没有关系。例如，把 \(\mathbb Z/6\mathbb Z\) 当成整数模，立即能看见非零元素也可能被非零整数零化；把一个算子看成 \(k[t]\) 的作用，便能把多项式恒等式翻译成模的关系。一个新定义若暂时显得空泛，可以先回到这两个例子，找出它所允许或排除的现象。对张量积、商模或等价函子等构造，则应实际写出映射，检查其良定义和复合关系。

\ProseParagraphBreak
本书给出章末习题的解答，但解答最好留到自己尝试之后再读。算错一个矩阵乘法，或发现一个映射不保持标量作用，常比顺畅地读过一页定义更能说明问题在哪里。代数的抽象有时像把许多不同的计算收进同一个句子；只有再把这个句子带回整数、多项式和矩阵，才能知道自己确实读懂了它。

\bigskip
\begin{flushright}
R. EverStarine

2026年9月24日
\end{flushright}
